\chapter{A Stored-Program Computer from the Fact}
\label{chap:stored-program-computer}

The first five chapters built the parts: a number, a duration, a notion of computation, a calculus in
embryo, and the informational skeleton of a field. This chapter assembles them into a single machine
--- a stored-program computer, grown entirely from the Fact, that holds a procedure as data and runs
it. The word machine is meant in the mathematician's sense, not the engineer's: there is no device
here, no dial and no readout, only a finite construction that happens to compute. What makes it a
\emph{stored-program} computer is the same move that made the von Neumann architecture more than a
calculator --- the program it runs is built from the very same Facts as the data it runs on, so that a
procedure is just another kind of record, able to be stored, inspected, compiled, and executed like
any other. Reaching this point completes the first part of the book: the construction has built itself
something to run on.

\section{The binary digit and its tick}
\label{se:the-binary-digit-and-its-tick}

A computer needs a unit of state that can be set, read, and changed, and the construction builds the
smallest one possible. It begins by noticing that two procedures can be \emph{equivalent} even when
they are written differently: if they carry every value to the same result, the construction treats
them as one, funging the procedures that act alike into a single equivalence class. Equivalence under
procedure is the chapter's first act --- the recognition that what a computation \emph{is} matters less
than what it \emph{does}, so that all the ways of doing the same thing collapse to one.

A small example shows why this funging is the right first move for a computer. Two procedures that add
one to a number --- one that counts up by a single tick, another that counts up by three and then back
by two --- are written differently, take different numbers of steps, and leave different intermediate
traces, yet they carry every input to the same output. For the purpose of building a machine they are
the same procedure, and the construction insists on treating them so. This is not laziness; it is the
precondition for compilation. A compiler may replace a procedure with any procedure equivalent to it
precisely because equivalence-under-procedure has already declared them interchangeable, and an
optimization that did not respect that equivalence would be changing the program's meaning rather than
its form.

On that footing the construction makes a binary digit. Recall from the first chapter that its numbers
carried not two truth-states but three: both-true, both-false, and the odd one out where the two
disagree. A clean binary digit has only two states, so the construction must \emph{exclude} one of the
three --- it tanges the two states that will count, zero and one, and sets the third aside. The digit
then carries a tick: a step that flips zero to one and one back to zero, the minimal nontrivial
dynamics a state can have. Set, read, flip: that is a bit, and the tick is what makes it a register
rather than a constant.

It is worth pausing on what was thrown away. The first chapter's numbers were richer than ordinary
bits: they carried three truth-states, not two, because they compared pairs of Facts and could register
agreement, disagreement, or the bare values themselves. A binary digit is poorer on purpose. The
construction tanges two of those three states to serve as zero and one and quarantines the third, and
the poverty is the point --- a clean two-state register is what a computer needs, and the construction
buys it by deliberately forgetting a distinction it was perfectly able to draw. Much of engineering is
this kind of disciplined forgetting: a usable component is an expressive structure with most of its
expressiveness suppressed, so that what remains is reliable.

The excluded third state does not vanish, and the construction gives it a name that the rest of the
book will earn. It calls the state that the binary digit leaves out --- the one the tick passes through
untouched, present but never selected --- a \emph{neutrino}, and the exclusion that produces it a
Pauli exclusion: only one of the two active states may be occupied at a time, exactly as only one
electron may occupy a state. This is the experiment under the neutrino's name, and it is careful to
disclaim any particle physics: the informational neutrino is simply the excluded state, the residue
that has to exist so that the digit can be cleanly binary, and it earns the analogy because it is the
thing that travels through the construction almost without interacting, carrying a small correction
that keeps the larger record globally consistent.

The choice of name is less arbitrary than it looks. A real neutrino is the particle that barely
interacts --- it streams through matter almost untouched, and from a supernova it arrives ahead of the
light because it need not wait for the medium to clear. The construction's excluded state plays the same
role in miniature: it passes through the tick unchanged while the two active states flip around it, and
it carries the small second-order correction that keeps separated observers' records agreeing with one
another. The analogy is doing real work --- the same informational shape, a nearly inert messenger that
reconciles the books --- and the construction claims only that shape, not the physics of weak
interactions or the mass of any particle.

The tick deserves one more word, because it answers an old objection. Berkeley complained that
Newton's fluxions were ghosts, increments that had to be nonzero to do their work and zero to
disappear afterward. The construction's tick is the fluxion made finite and honest: a single discrete
step, neither an infinitesimal nor a vanishing quantity, just the smallest change the digit admits.
This is the experiment under the name of fluxions, and its resolution is the one the whole book has
been rehearsing --- there are no infinitesimals here, and smooth change appears only as the unique
completion of finite distinctions, never as a quantity that is somehow both zero and not. A finite,
observable change emerges from finite, observable steps, and the ghost is laid to rest by refusing to
summon it.

\section{Source, compilation, and execution}
\label{se:source-compilation-and-execution}

With a bit and a tick the construction can build a program, and the way it does so is the heart of the
stored-program idea. A program begins as a source --- the construction's wry gloss is that the source
is wherever a procedure claims to have come from, an origin story that starts bending the moment it is
given a name. The source is then encoded into a sequence of operations, the opcodes, each a small,
fixed instruction the digital state knows how to obey. Compilation is the act of turning the source
into that encoded sequence; execution is the act of running it, each opcode ticking the state forward
one determined step.

The decisive point, and the reason this is a stored-program computer rather than a fixed device, is
that the encoded program is made of the same Facts as the data it manipulates. A program is not a
thing of a different kind, wired in from outside; it is a record, indistinguishable in substance from
the records it reads and writes. That single identity --- program is data --- is what lets a procedure
be stored alongside its inputs, inspected as if it were a number, transformed by another procedure, and
fed to itself. Everything that makes a modern computer general rather than special-purpose follows from
it.

The consequences are worth spelling out, because they are the difference between a calculator and a
computer. Because a program is a record, one program can read another, or read itself; it can rewrite a
procedure before running it, generate a new procedure as it goes, or hand a procedure to a further
procedure as input. A construction that kept programs in a separate, privileged compartment --- wired
in, unreadable, unchangeable --- could compute only the fixed function it was born with. By making the
program out of the same Facts as the data, the construction gets the whole of general-purpose
computation for free: anything it can do to a number it can do to a procedure, because in the end they
are the same kind of thing. The price of that generality is exactly the danger the next section
confronts, for a procedure that can manipulate procedures can manipulate itself.

A worked example shows the cycle turning. Take a source that says, in effect, \emph{flip the bit
three times}. Compilation encodes it as three identical tick-opcodes laid in order. Execution then runs
them: the state starts at zero, the first opcode ticks it to one, the second back to zero, the third to
one again, and the program halts with the bit reading one. Nothing in this required anything the
earlier chapters had not already built --- the bit is a Fact-built register, the opcodes are encoded
Facts, the order is the same order the number tower carried, and the halt is the computation of the
fourth chapter reaching its end. The computer is not a new primitive. It is the old primitives wired
together so that one record can direct the transformation of another.

Two features of this cycle reach back to earlier chapters and are worth naming. The compilation step is
where the source's claimed intent is fixed into a definite sequence of opcodes; after it, what runs is
the encoding, not the story the source told about itself, and any gap between the two is the program's
first chance to surprise its author. The execution step inherits the honesty of the fourth chapter: it
runs the opcodes one tick at a time and halts when it halts, with no guarantee issued in advance that it
will. A program that ticks forever is not forbidden by the construction; it is simply a procedure for
which the halting precondition was never established, and running it is the same act of faith every
computation here rests on. The computer does not escape the limits of computation by being assembled
out of Facts. It inherits them.

\section{Abstraction and the limit of self-reference}
\label{se:abstraction-and-self-reference}

A computer earns its power by abstraction: by hiding a procedure that works behind a name, so that a
solved problem can be reused as though it were a primitive. The construction builds this too, wrapping
compiled-and-executed procedures into abstractions that later procedures can call without seeing
inside. But abstraction is where a construction can hurt itself, because a procedure that can name and
manipulate procedures can be turned to name and manipulate \emph{itself}, and self-reference is the
doorway to paradox.

The construction guards the doorway with a move that is worth watching closely, because it is the last
structural decision of the part. It declines to make the order relation on abstractions total. Where
every earlier kind of number could be compared to every other, the order on abstractions is defined
only for certain pairs and deliberately left undefined for others --- in particular, for the pairs that
would let an abstraction ask a ranking question about itself. The motivation is the Berry paradox: the
smallest positive integer not definable in fewer than eleven words, a phrase of ten words that appears
to define exactly the integer it says is undefinable. The resolution is to deny that ``definable'' is
well defined across the board, and the construction makes the same move for ``less than'': by refusing
to rank certain self-referential pairs, it never lets the paradoxical question be posed, and the
machinery stays consistent.

Concretely, the construction can still compare two abstractions when their relationship is settled by
their structure --- a procedure that merely compiles ranks below one that compiles and then executes,
and such pairs are compared freely. What it refuses to do is rank an abstraction against a question
about its own ranking. The moment a comparison would require an abstraction to evaluate a predicate that
mentions that very comparison, the order returns no verdict at all, and the recursion that would
otherwise spiral into the Berry paradox has nowhere to begin. The fence is not a patch bolted on after
the fact; it is built into the definition of the order, which is total exactly where comparison is safe
and silent exactly where it is not.

This is the experiment the construction files under the paradoxes of time travel, and its lesson is
exact. The apparent paradoxes --- a message to one's own past, a procedure that ranks itself --- are
not glimpses of exotic physics but pathologies of over-resolution: they arise only when incompatible
refinements are treated as simultaneously admissible, when an observer tries to handle its own index as
an ordinary piece of data. The construction follows the warning that Godel made famous: a system that
can encode statements about its own inferences cannot, in general, stay globally consistent, so the
honest course is to keep certain self-referential configurations outside the admissible structure
altogether. The partial order is exactly that fence. It costs a little --- not everything can be
compared --- and it buys the one thing a computer built to reason must not spend: its own consistency.

\section{What is demonstrated, and what is assumed}
\label{se:demonstrated-assumed-ch6}

Following the experimental library's discipline one last time for this part, we separate the built from
the named, and note first what the chapter pointedly did \emph{not} use: no instrument, no dial, no
readout of any kind. The computer here is a finite mathematical object, and every step of it was grown
from Facts.

What is demonstrated is a working stored-program computer in miniature. There is equivalence under
procedure, funging procedures that act alike. There is a binary digit, tanged from the three truth-
states by excluding one, with a finite tick for its dynamics. There is a program built of the same
Facts as its data --- a source, an encoding into opcodes, a compilation, an execution that ticks the
state to a halt. And there is an abstraction mechanism made safe by a deliberately partial order that
keeps self-reference out. All of it is finite, decidable, and free of any physical device.

\begin{quote}
\emph{A note on the labels.} The physical names in this chapter are bridges, as before. Calling the
excluded state a \emph{neutrino}, the exclusion a \emph{Pauli exclusion}, and the finite tick a
\emph{fluxion} aligns the construction with the language of physics and the calculus, and each is
earned only to the extent its analogy holds: the construction exhibits the informational shape of the
excluded state and the discrete increment, and claims nothing about supernovae or about Newton's
dynamics. The one genuinely load-bearing identification is internal, not physical --- that a stored
program is the same kind of object as its data --- and that one is not an analogy but a fact about the
construction.

\end{quote}

What is assumed is the interpretive thread the whole part has been spending: that these finite
constructions deserve the physical and computational names laid over them. With the computer built, the
account has closed a loop --- the construction that began by measuring whether a single Fact is true has
assembled, out of nothing but such Facts, a machine capable of running the very procedures that decide
them. The chapters that follow leave the foundations behind and turn to what the machine is for: the
calculus of variations, the stationary actions, and the small finite field theory that the opening
chapters promised were waiting at the end of the construction.
